\documentclass{article}
\usepackage[a4paper,margin=18mm]{geometry}
\usepackage{unicode-math}
\setmathfont{XITS Math}
\usepackage{xcolor}
\pagestyle{empty}
\newbox\leftcontent
\newbox\rightcontent
\newbox\fullcontent
\newdimen\mathaxis
\newdimen\targetradius
\newcommand\LeftContent{a+b}
\newcommand\RightContent{(c+d)^2}
\newcommand\Body{\LeftContent+\frac{\displaystyle\sum_{k=1}^{n}x_k}{1+\frac pq}+\RightContent}
\newcommand\LeftSized{\mathopen{\nulldelimiterspace=0pt\left[\vphantom{\vrule width0pt height\dimexpr\targetradius+\mathaxis\relax depth\dimexpr\targetradius-\mathaxis\relax}\right.}}
\newcommand\RightSized{\mathclose{\nulldelimiterspace=0pt\left.\vphantom{\vrule width0pt height\dimexpr\targetradius+\mathaxis\relax depth\dimexpr\targetradius-\mathaxis\relax}\right]}}
\newcommand\Comparison[3]{\par\medskip\noindent\textbf{#1}\hfill Target radius: #3\,pt\par
 \targetradius=#2sp\relax
 \[ F=\mathinner{{\color{blue}\LeftSized}\copy\fullcontent{\color{blue}\RightSized}} \]
}
\begin{document}
\section*{Single-line delimiter sizing}
Identical content throughout: low endpoints and a tall middle fraction.
Sizes are blended as distances from the math axis, keeping both delimiters
centered on that axis. Native font size selection is retained.
\setbox\leftcontent=\hbox{$\displaystyle\LeftContent$}
\setbox\rightcontent=\hbox{$\displaystyle\RightContent$}
\setbox\fullcontent=\hbox{$\displaystyle\Body$}
\mathaxis=\Umathaxis\displaystyle
\directlua{dofile('experimental/single-line-size.lua')}
\medskip\noindent
Here $E$ is the larger endpoint radius; $H$ is the full-content radius.
The target radius is a sizing input: the font offers discrete delimiter sizes,
so different inputs can produce the same visible size.
\end{document}
